
\lettrine{E}{nfin, une trêve est établie. Les commissaires des deux partis passent pour donner l'ordre de démolir « les fortifications et marques de guerre... et temples de Trébas, Alban et quelques maisons rivales ». C'est la première fois que l'histoire signale un temple et sa communauté protestante en nos murs.}

\lettrine{C}{'est une paix bien précaire que celle établie en 1581. Les bagarres avaient repris en 1584 et Alban de nouveau retombait en mains protestantes. Aussi le Sénéchal Cornusson amène-t-il ses troupes de Toulouse, s'empare du château d'Alban. Ceci signifie une fois de plus : pillage. Au printemps suivant, dès que les troupes de Cornusson se sont retirées, les Protestants reviennent et, après avoir provoqué quelques incendies, se retirent par crainte d'un retour des soldats du sénéchal. Mais l'été de 1586 est à peine commencé qu'ils sont à nouveau là. Cette fois-ci, ils sont en force, ayant emmené avec eux deux pièces de batterie. Sont-elles semblables à celle qui se trouvent encore au Château de Paulin ? Ne serait-ce pas ces deux couleuvrines qu'on y voyait naguère ?... Quoi qu'il en soit, le clocher de l'église est transformé en forteresse provisoire. Les capitaines Comte et Caussanel y ont installé leurs cinquante hommes. Les Réformés peuvent attendre de pied ferme les troupes catholiques : elles seront chaudement reçues. Cela ne va pas tarder. Les troupes de l'évêque d'Albi ont rejoint celles du sénéchal et, après avoir attaqué Lombert, donnent l'assaut à Alban le 16 mai 1586.}

\lettrine{Q}{uelques jours avant l'assaut donné par les troupes catholiques, le capitaine Marchet s'était déjà emparé de quelques maisons hors ville. Le clocher, centre de la résistance protestante, est incendié. Vingt hommes y périssent avec le capitaine Comte. Le succès est de courte durée et une fois de plus, Alban tombe au pouvoir des réformés. Il y a eu cependant une contre offensive de la part des catholiques. Ils ont essayé d'occuper la place à la faveur d'un déguisement. « Le 28 du mois de may, dit un chroniqueur protestant, les papistes, accoutrés en femmes, ont pensé surprendre la ville d'Alban, de jour. Mais par la grâce de Dieu, un capitaine de la religion (=protestant) releva le pont promptement ». (Alban a donc encore son enceinte de murailles avec porte haute et porte basse nanties}
